\clearpage
\fancyhead[L]{\small MERKURTAFELN.}
\SourceLeaf{564}
\begin{center}\small Göttingische gelehrte Anzeigen. Stück 107. Seite 1065..1069. 5. Juli 1813.\\[3pt]\rule{.72\linewidth}{.3pt}\end{center}

Herr \textsc{von Lindenau}, korrespondierendes Mitglied der Königlichen Sozietät der Wissenschaften, hat sich mit seinen Venus- und Marstafeln bereits um die Planetentheorie verdient gemacht. In einem Schreiben vom 20. Juni dieses Jahres an Professor \textsc{Gauß} gab er eine vorläufige Nachricht über eine ähnliche, nicht weniger wichtige dritte Arbeit über den \emph{Merkur}, die er vor Kurzem abgeschlossen hat. Die verdienstlichen Bemühungen von \textsc{Lalande}, \textsc{Oriani} und \textsc{Triesnecker} hatten die Theorie der Merkurbahn zwar bereits zu einem hohen Grad an Vollkommenheit gebracht. Dennoch gab es praktische und theoretische Hilfsmittel, die teils noch nicht benutzt worden, teils neu hinzugekommen waren. Daher war es nicht unwahrscheinlich, dass eine neue Bearbeitung dieses Gegenstands weitere Änderungen und Verbesserungen der vorhandenen Bestimmungen liefern würde, wenn sie sorgfältig alles berücksichtigte, was die heutige Astronomie zu diesem Zweck bietet. Diese Aussicht war zusammen mit der Hoffnung, durch einen genauen Vergleich der theoretisch bestimmten Störungen der elliptischen Merkurbahn mit den aus den Beobachtungen folgenden eine neue und zuverlässige Bestimmung der Venusmasse ableiten zu können, der Hauptgrund, der Herrn \textsc{von Lindenau} zur neuen Bearbeitung des Gegenstands veranlasste. Wir teilen hier nur einige Endergebnisse mit. Die Einzelheiten dieser Untersuchungen werden den Astronomen in einem Werk dargelegt, das in wenigen Wochen unter dem Titel \emph{Investigatio nova orbitae a Mercurio circa Solem descriptae: accedunt tabulae Planetae ex elementis recens correctis et ex theoria gravitatis ill. de la Place constructae, auctore} \textsc{Bernhardo de Lindenau} im Druck erscheint.

Entsprechend dem zuvor angedeuteten Zweck gliedert sich die Untersuchung in vier Abschnitte. In den beiden ersten werden die seit 1631 beobachteten Merkurdurchgänge erörtert und zur Bahnbestimmung verwendet, in den beiden letzten neuere geozentrische Beobachtungen. Auf die Merkurdurchgänge wurden zwei Verfahren angewandt. Erstens ließen sich daraus der Knoten, die Knotenbewegung und die Neigung bestimmen. Zweitens hat die Merkurbahn die Besonderheit, dass alle Durchgänge im selben Knoten nahezu an denselben Punkten der Bahn stattfinden. Dadurch ließen sich die durch die Durchgänge gegebenen heliozentrischen Längen in der Ekliptik auf mittlere Längen in der Bahn reduzieren, ohne dass die Differenzen zweier Längen durch andere Elemente als die mittlere jährliche Bewegung und die jährliche Änderung des Aphels beeinflusst und wesentlich verfälscht werden konnten. Zur Anwendung dieses Verfahrens muss man allerdings die Merkurbahn bereits näherungsweise kennen, wie es hier tatsächlich der Fall war. Für die erste Bestimmung, die vor allem eine genaue Kenntnis der geozentrischen Breite erforderte, konnten nur zwölf Durchgänge benutzt werden. Aus den insgesamt siebzehn Durchgängen, die der Verfasser erörtert und berechnet hat, wurden fünfzehn Kombinationen gebildet, die ihm für die numerische Entwicklung und Anwendung des zweiten Verfahrens am besten geeignet erschienen. Der Vergleich der so erhaltenen jährlichen Änderungen des Knotens und des Aphels mit den theoretischen Änderungen, deren Werte hauptsächlich von der Venusmasse abhängen, lieferte die beiden ersten Gleichungen zur Bestimmung dieser Masse.

Die aus den eben erwähnten siebzehn Durchgängen abgeleiteten heliozentrischen Längen des Merkur dienten zusammen mit hundert von \textsc{Maskelyne} und \textsc{Piazzi} beobachteten geozentrischen Längen dazu, die eigentlich elliptischen Elemente der Merkurbahn zu bestimmen. Da für die Neigung und den Knoten bereits Näherungswerte vorlagen, konnten diese Elemente, die nur geringen \SourceLeaf{565}Einfluss auf die Reduktion haben, zunächst ohne Bedenken unberücksichtigt bleiben. Die Differenzen zwischen den beobachteten und den berechneten Längen wurden als Funktion der Korrekturen der Epoche, der mittleren Bewegung, der Exzentrizität, des Aphels und der Venusmasse ausgedrückt. So entstanden 117 Bedingungsgleichungen, aus denen die Korrekturen der zugrunde gelegten Elemente nach der Methode der kleinsten Quadrate abgeleitet wurden. Daraus ergab sich auch die dritte Gleichung für den Wert der Venusmasse. Bei der Aufstellung der Bedingungsgleichungen aus den geozentrischen Längen waren Vorsicht und Genauigkeit nötig. Bei den größten Elongationen, bei denen Merkur hauptsächlich und fast ausschließlich beobachtet worden ist, wird der heliozentrische Fehler durch die Reduktion auf den geozentrischen Ort stark vermindert. Kleine Fehler in den beobachteten geozentrischen Längen ändern daher die gesuchten Korrekturen der Elemente stark. Der Verfasser glaubt, diesen Einfluss der Beobachtungsfehler durch die Anzahl der Beobachtungen ausgeschaltet und den Einfluss der Sonnenörter auf die berechneten Längen vermieden zu haben, indem er die Sonnentafeln jedes Mal anhand gleichzeitiger Beobachtungen korrigierte.

Mit diesen verbesserten Elementen wurden die beobachteten geozentrischen Breiten auf heliozentrische Breiten reduziert. Aus ihrem Vergleich mit den berechneten Breiten ergaben sich hundert Bedingungsgleichungen, die die Neigung und den Knoten lieferten. Der Wunsch, die Neigung für eine frühere Epoche zu bestimmen und damit aus ihrer beobachteten säkularen Änderung noch eine vierte Gleichung für die Venusmasse zu erhalten, scheiterte am Mangel geeigneter Beobachtungen.

Die Endergebnisse dieser Untersuchungen waren die folgenden.
Setzt man die Venusmasse, wie sie \textsc{Laplace} (Mécan. Céleste T.\,3. p.\,61) annimmt, $= 1$, so ist die verbesserte Masse, wie sie aus den erwähnten drei Gleichungen folgt, $= 1{,}0974 = \frac{1}{849132}$ der Sonnenmasse. Mit dieser Masse wurden alle periodischen und säkularen Störungen berechnet.

\begin{center}\fontsize{9}{11}\selectfont\setlength{\tabcolsep}{2pt}\renewcommand{\arraystretch}{1.18}\begin{tabular}{@{}l c r@{\;}r@{\;}r@{}}
Epoche 1750, Meridian von Seeberg & $.\;.\;.\;.\;.\;.$ & $253^\circ$ & $5'$ & $17''1$\\
Mittlere jährliche Bewegung & $.\;.\;.\;.\;.\;.\;.\;.$ & $53$ & $43$ & $3.613$\\
Aphel 1750 & $.\;.\;.\;.\;.\;.\;.\;.\;.\;.\;.\;.\;.$ & $253$ & $33$ & $24.3$\\
Exzentrizität 1800 & $.\;.\;.\;.\;.\;.\;.\;.\;.\;.\;.\;.$ & \multicolumn{3}{l}{$0.2056163$}\\
Große Halbachse & $.\;.\;.\;.\;.\;.\;.\;.\;.\;.\;.\;.$ & \multicolumn{3}{l}{$0.3870988$}\\
Knoten 1750 & $.\;.\;.\;.\;.\;.\;.\;.\;.\;.\;.\;.\;.\;.\;.$ & $75^\circ$ & $22'$ & $0''96$\\
Neigung der Bahn 1800 & $.\;.\;.\;.\;.\;.\;.\;.\;.\;.$ & $7$ & $0$ & $5.9$\\
Säkulare Änderung des Aphels & $.\;.\;.\;.\;.\;.\;.$ & $+1$ & $33$ & $22.9$\\
\quad —\quad —\quad des Knotens & $.\;.\;.\;.\;.\;.\;.\;.$ & $+1$ & $10$ & $15.1$\\
\quad —\quad —\quad der Exzentrizität & $.\;.\;.\;.\;.\;.$ & \multicolumn{3}{l}{$+0''791$}\\
\quad —\quad —\quad der Neigung & $.\;.\;.\;.\;.\;.\;.\;.$ & \multicolumn{3}{l}{$+18.380$}\\
\end{tabular}\end{center}

{\rightskip=0pt plus 3em Diese Elemente stellen alle hundert beobachteten geozentrischen Örter, die in die Bedingungsgleichungen aufgenommen wurden, äußerst zufriedenstellend dar. Dasselbe gilt für fünfzig neuere Beobachtungen des Merkur, die \textsc{Bouvard} in Paris und der Verfasser selbst auf der Seeberger Sternwarte anstellten und die nicht zur Bestimmung der Elemente herangezogen wurden. Man ist daher durchaus zu der Hoffnung berechtigt, dass diese Bestimmungen die Bewegungen des Merkur auch in den nächsten Jahrzehnten in Übereinstimmung mit dem Himmel darstellen werden.\par}
